\section{历史基础设施周期：稀缺、过建与标准化}

讲者把今天的 compute 放进更长的基础设施史：钢铁、铁路、电力、光纤、半导体材料与能源技术都经历过需求爆发、资本涌入、供给扩张、价格下降、企业整合和标准形成。历史类比的用途是生成问题，比如“哪里会过建”“什么接口会标准化”“谁承担沉没成本”；它不能直接预测时间表，因为每种资源的物理性质、监管结构和创新速度不同。

\subsection{Gilded Age：基础设施与权力一起扩张}

19 世纪末的 Gilded Age 常被用来描述铁路、钢铁和金融快速扩张，同时市场力量集中、劳动冲突和公共监管加剧的阶段。讲者借它提醒学生，通用技术的繁荣不仅带来更多产能，也会重新分配议价权。AI 基础设施的关键问题因此不止是“能建多少”，还包括谁控制接入、谁制定价格、谁能在短缺时继续创新。

\begin{figure}[H]
\centering
\includegraphics[width=0.88\textwidth]{images/00-51-00-gilded-age.jpg}
\caption{Gilded Age 为本讲提供历史入口：基础设施繁荣与权力集中往往同时出现。}
\figsource{Stanford Online 官方视频 00:51:00，`images/00-51-00-gilded-age.jpg`。}
\end{figure}

\begin{knowledgebox}{读图：历史类比先比较机制}
可比较的机制包括高固定成本、网络效应、资源控制、纵向整合、价格歧视与公共监管压力。不可直接类比的部分包括技术更新速度、全球供应链、软件复制成本和现代资本市场。把“今天就是新的 Gilded Age”当结论过于粗糙；把它当成检查市场集中与公共利益的提示，更有教学价值。
\end{knowledgebox}

\subsection{钢铁：工艺改进怎样改变资源价格}

Bessemer process 等工艺创新显著降低钢铁生产成本，使铁路、建筑和机械能够在更大范围扩张。这个案例对应 AI 中的模型效率、芯片设计和系统优化：需求增长不一定只靠增加原材料供给来满足，生产过程本身也会改变单位成本。价格下降又会释放新需求，使总消费量继续上升。

\begin{figure}[H]
\centering
\includegraphics[width=0.94\textwidth]{images/00-52-00-bessemer-steel-prices.jpg}
\caption{Bessemer 工艺之后钢铁价格长期下降，生产效率改变了基础设施可及性。}
\figsource{Stanford Online 官方视频 00:52:00，`images/00-52-00-bessemer-steel-prices.jpg`。}
\end{figure}

\begin{importantbox}{读图：单位成本下降可能扩大总需求}
钢铁更便宜以后，社会没有停止使用钢铁，而是把它带入更多铁路、桥梁和建筑。计算效率也可能产生类似的 rebound effect：每次推理更便宜，会催生更多 agent、更长任务和更高频交互。判断能源与算力总需求时，不能只看单次任务成本下降，还要估计新用途带来的需求弹性。
\end{importantbox}

\subsection{光纤：过度建设也可能留下公共能力}

互联网泡沫时期铺设的大量光纤在短期造成破产与资产折价，却为后续宽带和云服务留下低成本容量。这个故事常被用来为 AI 数据中心扩张辩护，但两者仍有关键差异：光纤暗纤可以长期闲置，GPU 折旧更快，芯片代际更新、封装和软件兼容都会缩短有效寿命。过建是否有长期公共价值，取决于资产能否被重新利用。

\begin{figure}[H]
\centering
\includegraphics[width=0.94\textwidth]{images/00-53-00-fiber-optic-overbuild.jpg}
\caption{光纤价格在过度建设后快速下降，闲置容量后来成为互联网扩张的基础。}
\figsource{Stanford Online 官方视频 00:53:00，`images/00-53-00-fiber-optic-overbuild.jpg`。}
\end{figure}

\begin{warningbox}{读图：``先过建再说'' 不是无条件策略}
资产寿命、维护成本、标准兼容和再分配难度决定过建能否转化为公共能力。数据中心若绑定单一电网、封闭互连或快速过时的芯片，剩余价值可能很低。基础设施规划应当同时考虑 modularity、可升级性和开放接口，而不是只比较峰值建设规模。
\end{warningbox}

\subsection{DRAM：半导体供给为何呈周期性}

\term{DRAM} 是 Dynamic Random-Access Memory，中文常译为动态随机存取存储器。它以电容保存 bit，需要周期刷新，广泛用作服务器和个人计算机主存。晶圆厂投资周期长、固定成本高，而需求又会随设备与数据中心景气变化，因此价格常在短缺与过剩之间摆动。AI 加速器所用显存与先进封装也受到类似产能规划约束。

\begin{figure}[H]
\centering
\includegraphics[width=0.94\textwidth]{images/00-53-30-dram-cycles.jpg}
\caption{DRAM 产业周期展示扩产滞后、库存调整和价格波动如何反复出现。}
\figsource{Stanford Online 官方视频 00:53:30，`images/00-53-30-dram-cycles.jpg`。}
\end{figure}

\begin{knowledgebox}{读图：供给滞后怎样制造周期}
需求上升时价格先涨，厂商据此扩产；新产能经过建设和爬坡后集中进入，市场可能转为过剩，价格下跌又迫使企业削减投资。等需求恢复时，上一轮削减导致供给再次紧张。GPU、HBM 和先进封装并不与 DRAM 完全相同，但都需要把长建设周期和短需求预测放在同一模型中分析。
\end{knowledgebox}

\subsection{铀与能源：政策、许可和地缘因素改变曲线}

算力扩张最终落到能源，因此本节从芯片与机房继续追到发电、燃料和电网。铀市场说明，资源价格不仅受矿山供给与需求影响，也受核电政策、许可周期、库存、地缘政治和公众接受度影响。即使一种能源在技术上适合稳定供电，它能否及时进入数据中心仍取决于电网互连、融资、建设和监管协调。读图时要把矿产价格、可投运电力与特定园区真正可获得的容量分开，三者之间存在多年建设与审批滞后。

\begin{figure}[H]
\centering
\includegraphics[width=0.94\textwidth]{images/00-54-15-uranium-cycle.jpg}
\caption{铀价格周期反映资源、政策、库存与能源需求共同塑造基础设施成本。}
\figsource{Stanford Online 官方视频 00:54:15，`images/00-54-15-uranium-cycle.jpg`。}
\end{figure}

\begin{knowledgebox}{读图：价格不是纯粹的物理稀缺度}
矿藏数量只决定供给的一部分，开采许可、转化与浓缩能力、长期合同、战略库存和核电建设计划都会改变可用供给。AI 能源讨论也要避免把 GW 数字当作随时可调用的电力：地点、时间、稳定性、输电容量和碳约束决定它能否服务具体集群。
\end{knowledgebox}

\subsection{Alternative assets：同一种波形背后可能是不同机制}

讲者把多类 alternative assets 的价格周期放在一起，试图寻找从稀缺、追捧、扩产到回落的共同时间结构。这种比较可以帮助投资与建设团队建立情景范围，却也容易产生“所有周期都一样”的错觉。曲线相似可能来自不同原因，持续时间也受资产寿命、监管和替代品影响。

\begin{figure}[H]
\centering
\includegraphics[width=0.94\textwidth]{images/00-55-45-alternatives-timeline.jpg}
\caption{多种 alternative assets 的周期对照，用于讨论稀缺叙事、资本涌入和供给响应。}
\figsource{Stanford Online 官方视频 00:55:45，`images/00-55-45-alternatives-timeline.jpg`。}
\end{figure}

\begin{warningbox}{读图：形状相似不等于因果相同}
价格先涨后跌可能来自扩产、需求崩塌、政策变化、技术替代或金融去杠杆。若只做视觉匹配，就会忽略每种资产的交付周期、质量差异和网络效应。使用历史类比时，应先写出机制，再检查当前 compute 市场是否满足同样条件。
\end{warningbox}

\begin{table}[H]
\centering
\begin{tabular}{L{0.14\textwidth}L{0.21\textwidth}L{0.23\textwidth}L{0.25\textwidth}}
\toprule
案例 & 可迁移机制 & 对 compute 的启发 & 不能直接类比之处 \\
\midrule
钢铁 & 工艺创新降低单位成本 & 软件与芯片效率改变可及性 & 数字需求增长更快 \\
光纤 & 过建后资产折价再利用 & 预建容量可能支持后续应用 & GPU 折旧和代际更快 \\
DRAM & 长扩产周期与短需求周期错配 & 芯片、封装与库存会波动 & 产品异构和供应链更复杂 \\
铀 & 政策、许可与地缘因素塑造供给 & 电力不是抽象无限资源 & 能源与计算的交付接口不同 \\
Gilded Age & 网络基础设施带来市场集中 & 接入与公共利益需要制度关注 & 当代全球监管结构不同 \\
\bottomrule
\end{tabular}
\caption{历史类比的正确用法：迁移机制，同时保留物理与制度差异。}
\end{table}

\subsection{本章小结}

历史周期提供了五类提醒：工艺创新会降低单位成本，需求反弹会扩大总用量，建设滞后会制造价格周期，过建价值取决于资产可复用性，市场集中会推动标准与公共监管。下一节把这些提醒用于一个具体判断：今天的 compute 还不能像电力那样被当作可替代商品。

